\documentclass[12pt,a4paper]{article}
\usepackage[utf8]{inputenc}
\usepackage[T1]{fontenc}
\usepackage{geometry}
\usepackage{enumitem}
\usepackage{hyperref}
\usepackage{parskip}
\geometry{margin=1in}

\begin{document}

\begin{titlepage}
    \centering
    \vspace*{2cm}
    {\LARGE \textbf{SocsBoard} \\[1.5em]}
    {\Large Software Requirements Specification \\[1em]}
    {\large Version 1 \\[0.5em] 7/11/25 \\[2em]}
    {\large Caoilfhionn Fitzpatrick, Samuel Regan, Semyon Fox, Shreyansh Singh \\[4em]}
    Prepared for\\
    CT216---Software Engineering I
    \vfill
\end{titlepage}

\tableofcontents
\newpage

\section{Introduction}

\subsection{Purpose}

This Software Requirements Specification (SRS) document provides a comprehensive description of the University Society Platform (SocsBoard) project requirements. It details the functional and non-functional requirements, user stories, and acceptance criteria for the platform.

\textbf{The SRS is intended for:}
\begin{itemize}
  \item The development team
  \item The project manager (Semyon Fox) for reviews and approval management
  \item Module instructors and academic assessors
  \item Future maintainers and contributors
  \item Stakeholders interested in the technical scope and deliverables
\end{itemize}

\textbf{This document serves as the foundation for:}
\begin{itemize}
    \item Development planning and task allocation
    \item Testing and quality assurance
    \item Project progress tracking
    \item Academic evaluation
\end{itemize}

\subsection{Key Stakeholders}

\begin{enumerate}[label=(\arabic*)]
    \item \textbf{Project Manager:} Semyon
    \begin{enumerate}
        \item Responsible for pull request reviews and approvals
        \item Manage merge conflicts and code quality
        \item Facilitates architecture decisions
        \item Coordinates team tasks
    \end{enumerate}

    \item \textbf{Development Team:} 4 full-stack developers
    \begin{enumerate}
        \item Build features through, end to end (frontend, backend, database)
        \item Write tests and documentation for their features
        \item Participate in code reviews
        \item Own 2-3 features each
    \end{enumerate}

    \item \textbf{End Users:}
    \begin{enumerate}
        \item \textbf{Students:} Discover society events and updates, view society posts/events, show interest in an event and receive personalised recommendations
        \item \textbf{Society Committee:} Create and manage events, posted content and view engagement analytics
        \item \textbf{Sudo:} Platform management, user moderation, system monitoring
    \end{enumerate}

    \item \textbf{University Administration:}
    \begin{enumerate}
        \item Provide OAuth authentication services
        \item Supply student data (by API access)
        \item Approve production deployment
    \end{enumerate}

    \item \textbf{Academic Assessors:}
    \begin{enumerate}
        \item Module teaching staff evaluating deliverables
        \item Assess technical implementation, documentation and knowledge
    \end{enumerate}
\end{enumerate}

\subsection{Scope, Objectives and Goals}

\subsubsection{Software System Identified}

The \textbf{SocsBoard (University Society Platform)} is a full-stack web-based social and event platform with a recommendation engine that connects students with society events and social content through personalized relevance scoring.

\textbf{Current Deployment:} The platform is currently hosted at \url{https://ct216.semyon.ie} using a Cloudflare Tunnel. The application runs securely within Docker containers with all databases isolated in a dedicated subnet on a private server infrastructure.

\subsubsection{Scope}

\textbf{What the system WILL do:}
\begin{itemize}
    \item Provide full event management capabilities, including creation, promotion, capacity tracking, and multilingual descriptions
    \item Support social content sharing through posts, images, and informal updates
    \item Deliver personalized recommendations based on user behavior, interests, and society memberships
    \item Handle authentication through university OAuth (Microsoft Azure AD) with automatic synchronization of student data (course, department, year, memberships) (Subject to approval by ISS)
    \item Ensure GDPR compliance and support both English and Irish UI languages, with optional translation for user-generated content
    \item Allow users to create events in English, provide their own translations, or use AI for automatic translations
    \item Enable a separate translator role to manage translations via a dedicated dashboard
    \item Support ticket-based human translation requests for educational purposes, such as students learning a language
    \item Provide analytics dashboards displaying engagement trends, demographics, and event performance metrics
\end{itemize}

\textbf{What the system will NOT do:}
\begin{itemize}
    \item Payment processing, ticketing, or any financial transactions
    \item Private messaging or chat functionality
    \item Academic features (grades, timetables, assignments)
\end{itemize}

\subsubsection{Benefits}

\textbf{\textit{For Students:}}
\begin{itemize}
    \item Discover relevant events through personalized recommendations
    \item Access all society activity in one centralized location
    \item Benefit from multilingual support
    \item Easily register for and track event attendance
\end{itemize}

\textbf{\textit{For Societies:}}
\begin{itemize}
    \item Promote events centrally
    \item Share both formal and casual content
    \item Analyze engagement metrics
    \item Reach students beyond their current membership
\end{itemize}

\textbf{\textit{For the University:}}
\begin{itemize}
    \item Gain improved engagement insights
    \item Track participation rates
    \item Enhance student experience
    \item Obtain data-driven visibility into campus life
\end{itemize}

\subsubsection{Technical Objectives}

\textbf{Learning and Development:}
\begin{itemize}
    \item Develop full-stack expertise using Next.js, React, PostgreSQL, and Redis
    \item Learn server-side rendering, client-side rendering, caching strategies, and OAuth implementation
    \item Adopt DevOps practices including Docker, CI/CD pipelines, and AWS deployment
    \item Build a multilingual application from the ground up
\end{itemize}

\textbf{Performance Targets:}
\begin{itemize}
    \item Fast page-load times for slower mobile connections, or bad WiFi coverage areas on campus
    \item Support for 100+ concurrent users
    \item Efficient Redis caching implementation
\end{itemize}

\textbf{Security Requirements:}
\begin{itemize}
    \item Secure OAuth implementation with no password storage
    \item Protection against XSS/CSRF/SQL-injection attacks
    \item GDPR compliance
\end{itemize}

\textbf{Scalability Goals:}
\begin{itemize}
    \item Database architecture that supports potentially thousands of users
    \item Extensible system design
    \item Cache-based recommendation engine
    \item Horizontal scaling capability
\end{itemize}

\subsubsection{Project Goals}

\textbf{\textit{Academic:}}
\begin{itemize}
    \item Deliver a functional prototype by Christmas with documentation and demonstrations
    \item Achieve high module performance
    \item Demonstrate strong software engineering practices throughout development
\end{itemize}

\textbf{\textit{Career Development:}}
\begin{itemize}
    \item Produce a portfolio-quality project showcasing industry-standard tooling experience
    \item Improve teamwork and code review skills
    \item Create compelling interview talking points
\end{itemize}

\textbf{\textit{Real-World Impact:}}
\begin{itemize}
    \item Address the genuine problem of student event discovery
    \item Enable potential campus deployment
    \item Gather feedback from actual users
    \item Iterate based on live usage patterns
\end{itemize}

\subsection{References}

\subsubsection{Project Documentation}

This SRS references the following internal project documents:

\begin{enumerate}
    \item \textbf{PROJECT\_OVERVIEW.md}
    \begin{itemize}
        \item Complete project scope and evolution
        \item 8-week timeline
        \item Team structure and approach
        \item Source: \texttt{/docs/PROJECT\_OVERVIEW.md}
    \end{itemize}

    \item \textbf{ARCHITECTURE\_AND\_TECH\_STACK.md}
    \begin{itemize}
        \item Technical architecture decisions
        \item Database schema design
        \item Technology stack rationale
        \item Source: \texttt{/docs/ARCHITECTURE\_AND\_TECH\_STACK.md}
    \end{itemize}

    \item \textbf{INFRASTRUCTURE\_AND\_DEVOPS.md}
    \begin{itemize}
        \item Docker setup and deployment
        \item CI/CD pipeline configuration
        \item AWS migration strategy
        \item Source: \texttt{/docs/INFRASTRUCTURE\_AND\_DEVOPS.md}
    \end{itemize}

    \item \textbf{REDIS\_CACHING\_GUIDE.md}
    \begin{itemize}
        \item Redis implementation details
        \item Caching strategies and patterns
        \item Performance optimization
        \item Source: \texttt{/docs/REDIS\_CACHING\_GUIDE.md}
    \end{itemize}

    \item \textbf{Stack Diagram.excalidraw.md}
    \begin{itemize}
        \item Visual system architecture
        \item Component relationships
        \item Source: \texttt{/docs/Stack Diagram.excalidraw.md}
    \end{itemize}
\end{enumerate}

\subsubsection{External References}

The following external documentation and technical resources are referenced:

\begin{itemize}
    \item \textbf{Next.js Documentation} (v14+) \\
    Available at: \url{https://nextjs.org/docs}

    \item \textbf{React Documentation} (v18+) \\
    Available at: \url{https://react.dev}

    \item \textbf{PostgreSQL Documentation} (v15+) \\
    Available at: \url{https://www.postgresql.org/docs/15/}

    \item \textbf{Redis Documentation} (v7+) \\
    Available at: \url{https://redis.io/docs}

    \item \textbf{Bootstrap Documentation} (v5.3) \\
    Available at: \url{https://getbootstrap.com/docs/5.3/}

    \item \textbf{next-intl Documentation} \\
    Available at: \url{https://next-intl-docs.vercel.app}

    \item \textbf{Microsoft Azure AD OAuth Documentation} \\
    Available at: \url{https://learn.microsoft.com/en-us/azure/active-directory/}
\end{itemize}

\subsubsection{Module Requirements}

\begin{itemize}
    \item \textbf{CT216 Software Engineering I Course Outline} \\
    University College, Academic Year 2025/26

    \item \textbf{Project Brief:} Full-stack web application using HTML, CSS, JavaScript, and Node.js

    \item \textbf{Assessment Criteria:} 40\% of module grade

    \item \textbf{Submission Deadlines:}
    \begin{itemize}
        \item Prototype: Christmas break
        \item Total project time: 6 months
    \end{itemize}
\end{itemize}

\subsection{Overview}

This SRS document is organized into the following sections:

\textbf{Section 1: Introduction} - Establishes the document's purpose, identifies key stakeholders and their roles, defines the project scope with clear boundaries, lists reference documents, and provides this structural overview.

\textbf{Section 2: Functional Requirements} - Contains user stories organized by epic (Authentication, Events, Social Posts, Recommendations, etc.) with acceptance criteria, plus interface requirements covering UI design philosophy and API structure.

\textbf{Section 3: Project Timeline} - Outlines development phases, key milestones, and risk management strategies. Detailed Gantt charts maintained separately.

\textbf{Appendices} - Supplementary material including glossary and additional technical details. Appendix content is informational and not part of the formal requirements unless explicitly stated.

\newpage
\section{Functional Requirements}

\subsection{User Stories}

\subsubsection{Epic 1: User Authentication \& Profile Management}

\textbf{User Story 1.1: University OAuth Login}
\begin{itemize}
    \item ``As a student or society member, I want to log in using my university account credentials so that I can access the platform securely without creating a separate password.''
    \item \textbf{Acceptance Criteria:}
    \begin{enumerate}
        \item User can click ``Login with University Account'' button on the landing page
        \item User is redirected to Microsoft Azure AD authentication page
        \item System receives OAuth callback and exchanges authorization code for access token
        \item System creates or updates user record in the database
        \item User is redirected to personalized dashboard
        \item User remains logged in across browser sessions
    \end{enumerate}
\end{itemize}

\textbf{User Story 1.2: User Profile Setup}
\begin{itemize}
    \item ``As a new user, I want to complete my profile with interests and preferences so that I receive personalized recommendations.''
    \item \textbf{Acceptance Criteria:}
    \begin{enumerate}
        \item After first login, user is prompted to complete profile
        \item User can select interests from predefined categories (sports, music, technology, arts, etc.)
        \item User can set language preference (Irish or English)
        \item Changes are saved and reflected immediately
        \item Recommendations are updated based on new interests
    \end{enumerate}
\end{itemize}

\textbf{User Story 1.3: Session Management}
\begin{itemize}
    \item ``As a logged-in user, I want to stay logged in securely across sessions so that I don't have to re-authenticate every visit.''
    \item \textbf{Acceptance Criteria:}
    \begin{enumerate}
        \item User session is stored in Redis with TTL
        \item User can manually log out, invalidating their session
        \item User sees ``Session expired'' message after inactivity period
        \item Multiple concurrent sessions are supported
    \end{enumerate}
\end{itemize}

\subsubsection{Epic 2: Event Discovery \& Management}

\textbf{User Story 2.1: Browse Events}
\begin{itemize}
    \item ``As a student, I want to browse upcoming events so that I can discover activities that interest me.''
    \item \textbf{Acceptance Criteria:}
    \begin{enumerate}
        \item Landing page shows list of upcoming events sorted by date
        \item Each event card displays: title, society name, date/time, location, thumbnail
        \item Events show remaining capacity if limited
        \item User can filter by date range, society, and category
        \item User can search events by keyword
        \item ``Recommended for You'' section appears at top
    \end{enumerate}
\end{itemize}

\textbf{User Story 2.2: View Event Details}
\begin{itemize}
    \item ``As a student, I want to view detailed information about an event so that I can decide whether to attend.''
    \item \textbf{Acceptance Criteria:}
    \begin{enumerate}
        \item Event detail page shows full description, date, location, society info
        \item Shows number of registered attendees and capacity
        \item ``Register'' button visible (or ``Registered'' if already registered)
        \item Translation badge shows if content is auto-translated
    \end{enumerate}
\end{itemize}

\textbf{User Story 2.3: Register for Event}
\begin{itemize}
    \item ``As a student, I want to register for an event so that I can attend and receive updates.''
    \item \textbf{Acceptance Criteria:}
    \begin{enumerate}
        \item User clicks ``Register'' button on event detail page
        \item If capacity limit reached, user is added to waitlist
        \item User sees confirmation message
        \item User can unregister by clicking ``Cancel Registration''
        \item Capacity count updates accordingly
    \end{enumerate}
\end{itemize}

\textbf{User Story 2.4: Create Event (Society Admin)}
\begin{itemize}
    \item ``As a society committee member, I want to create a new event so that students can discover and attend it.''
    \item \textbf{Acceptance Criteria:}
    \begin{enumerate}
        \item User with society admin role sees ``Create Event'' button
        \item Form includes: title, description, date/time, location, capacity, image, category
        \item Language options: provide both languages, auto-translate, or single language
        \item Event can be saved as draft or published
        \item Event appears in listings immediately upon publishing
    \end{enumerate}
\end{itemize}

\textbf{User Story 2.5: Edit/Delete Event}
\begin{itemize}
    \item ``As a society committee member, I want to edit or delete an event so that I can correct mistakes or cancel events.''
    \item \textbf{Acceptance Criteria:}
    \begin{enumerate}
        \item Society admin sees ``Edit'' and ``Delete'' buttons on their events
        \item Edit opens form pre-filled with existing data
        \item Delete requires confirmation
        \item Registered users notified if event is cancelled
    \end{enumerate}
\end{itemize}

\textbf{User Story 2.6: Calendar Integration}
\begin{itemize}
    \item ``As a student, I want to add events to my personal calendar so that I don't miss events I'm interested in.''
    \item \textbf{Acceptance Criteria:}
    \begin{enumerate}
        \item ``Add to Calendar'' button visible on event detail pages
        \item Generates .ics file compatible with major calendar applications
        \item Calendar integration for registered events and expressed interest
        \item Calendar entries include: event title, date/time, location, description, society info
        \item Supports Google Calendar, Apple Calendar, Outlook, and other iCalendar-compatible applications
    \end{enumerate}
\end{itemize}

\subsubsection{Epic 3: Social Posts \& Society Interaction}

\textbf{User Story 3.1: Create Social Post}
\begin{itemize}
    \item ``As a society member, I want to post casual content so that I can engage with students informally.''
    \item \textbf{Acceptance Criteria:}
    \begin{enumerate}
        \item Society member sees ``Create Post'' button
        \item Form includes text content and optional image upload
        \item Post appears in society feed immediately
    \end{enumerate}
\end{itemize}

\textbf{User Story 3.2: View Social Feed}
\begin{itemize}
    \item ``As a student, I want to see posts from societies I follow so that I can stay updated.''
    \item \textbf{Acceptance Criteria:}
    \begin{enumerate}
        \item User sees ``Social Feed'' section on dashboard
        \item Feed shows posts from followed societies
        \item Posts display: society name, content, image, timestamp
        \item Posts sorted by recency or relevance
    \end{enumerate}
\end{itemize}

\textbf{User Story 3.3: Interact with Posts}
\begin{itemize}
    \item ``As a student, I want to like or ignore posts so that the system learns my preferences.''
    \item \textbf{Acceptance Criteria:}
    \begin{enumerate}
        \item User can click ``Like'' button on any post
        \item User can click ``Not Interested'' to hide post
        \item Interactions recorded for recommendation algorithm
    \end{enumerate}
\end{itemize}

\textbf{User Story 3.4: Follow/Unfollow Society}
\begin{itemize}
    \item ``As a student, I want to follow societies I'm interested in so that I see their events and posts.''
    \item \textbf{Acceptance Criteria:}
    \begin{enumerate}
        \item ``Follow'' button on society profile page
        \item Button toggles to ``Following'' when clicked
        \item User can unfollow at any time
        \item Followed societies appear on user dashboard
    \end{enumerate}
\end{itemize}

\subsubsection{Epic 4: Personalized Recommendations}

\textbf{User Story 4.1: Receive Event Recommendations}
\begin{itemize}
    \item ``As a student, I want to see personalized event recommendations so that I discover relevant activities without searching.''
    \item \textbf{Acceptance Criteria:}
    \begin{enumerate}
        \item ``Recommended for You'' section on dashboard
        \item Shows top recommended events
        \item Each recommendation shows reason (``Because you attended...'', ``Popular in your department'')
        \item Recommendations based on: society memberships, past attendance, stated interests, department
        \item User can dismiss recommendations (``Not Interested'')
    \end{enumerate}
\end{itemize}

\textbf{User Story 4.2: Recommendation Feedback}
\begin{itemize}
    \item ``As a student, I want to provide feedback on recommendations so that the system improves over time.''
    \item \textbf{Acceptance Criteria:}
    \begin{enumerate}
        \item Each recommendation has ``Interested'' and ``Not Interested'' buttons
        \item Feedback influences future recommendations
        \item User can see why content was recommended
    \end{enumerate}
\end{itemize}

\subsubsection{Epic 5: Multilingual Support}

\textbf{User Story 5.1: Switch Language}
\begin{itemize}
    \item ``As a user, I want to switch between English and Irish so that I can use the platform in my preferred language.''
    \item \textbf{Acceptance Criteria:}
    \begin{enumerate}
        \item Language selector visible in navigation header
        \item UI elements change immediately on selection
        \item Language preference saved
        \item User-generated content displays in selected language if available
        \item Fallback to English with notice if translation unavailable
    \end{enumerate}
\end{itemize}

\textbf{User Story 5.2: Create Multilingual Content}
\begin{itemize}
    \item ``As a society admin, I want to create content in both English and Irish so that all students can engage.''
    \item \textbf{Acceptance Criteria:}
    \begin{enumerate}
        \item Event/post form has tabs for both languages
        \item User can provide manual translations
        \item Auto-translate option available (DeepL API)
        \item Translation quality badge displayed on content
    \end{enumerate}
\end{itemize}

\textbf{User Story 5.3: Upload Multilingual Posters}
\begin{itemize}
    \item \textbf{Description:} As a society admin, I want to upload posters with optional translations so that events are visually appealing and accessible to all.
    \item \textbf{Acceptance Criteria:}
    \begin{enumerate}
        \item Event creation form includes an option to upload posters.
        \item Posters can have an optional Irish version.
        \item Posters are displayed alongside event details.
        \item Translator role can add additional language versions of posters.
    \end{enumerate}
\end{itemize}

\textbf{Future Scope:}
\begin{itemize}
    \item Expand language support beyond English and Irish to include other languages such as Hindi, Chinese, French, etc.
    \item Enable crowdsourced translations for user-generated content and posters.
    \item Build a community-driven translation system with incentives for contributors.
\end{itemize}

\subsubsection{Epic 6: Society Dashboard \& Analytics}

\textbf{User Story 6.1: View Society Dashboard}
\begin{itemize}
    \item ``As a society committee member, I want to access a dashboard with all management tools so that I can efficiently manage events and posts.''
    \item \textbf{Acceptance Criteria:}
    \begin{enumerate}
        \item Dashboard shows: upcoming events, recent posts, follower count
        \item Quick actions: ``Create Event'', ``Create Post''
        \item Analytics summary visible
    \end{enumerate}
\end{itemize}

\textbf{User Story 6.2: View Engagement Analytics}
\begin{itemize}
    \item ``As a society committee member, I want to see analytics about engagement so that I can understand what content resonates.''
    \item \textbf{Acceptance Criteria:}
    \begin{enumerate}
        \item Analytics show: registration trends, attendance rates, post engagement, follower growth
        \item Demographics breakdown: department, year of study
        \item Data visualized with charts
        \item Filterable by date range
    \end{enumerate}
\end{itemize}

\subsubsection{Epic 7: Administrative Functions}

\textbf{User Story 7.1: Moderate Content}
\begin{itemize}
    \item ``As a super admin, I want to review and moderate content so that the platform remains safe and appropriate.''
    \item \textbf{Acceptance Criteria:}
    \begin{enumerate}
        \item Admin panel shows reported posts and events
        \item Admin can delete or hide inappropriate content
        \item Admin can ban users with confirmation
        \item Actions logged in audit trail
    \end{enumerate}
\end{itemize}

\textbf{User Story 7.2: View Platform Analytics}
\begin{itemize}
    \item ``As a super admin, I want to view platform-wide analytics so that I can monitor system health and usage.''
    \item \textbf{Acceptance Criteria:}
    \begin{enumerate}
        \item Shows: total users, active users, total events, total posts
        \item Most popular societies
        \item Trends visualized with charts
    \end{enumerate}
\end{itemize}

\subsection{Interface Requirements}

\subsubsection{User Interfaces}

\textbf{Design Philosophy:}

\begin{table}[h]
\centering
\begin{tabular}{|p{4cm}|p{3cm}|p{7cm}|}
\hline
\textbf{Interface} & \textbf{Approach} & \textbf{Rationale} \\
\hline
Society Dashboard / Admin & Desktop-first & Complex data tables, analytics charts, multi-field forms - optimized for larger screens \\
\hline
Student Interface & Mobile-first & Casual browsing, quick event check-ins, social feed scrolling - primary use case is phone \\
\hline
\end{tabular}
\end{table}

\textit{Note: Future consideration for Cordova wrapper for app store distribution - trades native experience for increased app size.}

\textbf{UI/UX Design Principles:}

\begin{itemize}
    \item \textbf{Consistent \& Modern:} Clean, unified design language across all interfaces with smooth transitions and modern aesthetics
    \item \textbf{Accessible:} Easy to use for all users regardless of technical ability
    \item \textbf{Feed-Based Discovery:} Instagram-like scrolling experience for recommendations, designed to be engaging without being ``doomscrolly''
    \item \textbf{Perplexity-Inspired Sidebar:} Clean navigation with main options:
    \begin{itemize}
        \item Search
        \item Explore
        \item Societies
        \item Events
        \item Calendar
        \item Dashboard (conditional - only visible to society admins)
    \end{itemize}
\end{itemize}

\textbf{Unauthenticated Access:}
\begin{itemize}
    \item Platform viewable without login (non-personalized results)
    \item All user interactions blocked until authentication
    \item Browse events, societies, and content in read-only mode
    \item Login required for: following, liking, registering, posting, analytics
\end{itemize}

\textbf{Key Interface Components:}

\begin{enumerate}
    \item \textbf{Landing Page (Unauthenticated)}
    \begin{itemize}
        \item Hero section with platform tagline and login button
        \item Featured events preview (non-personalized)
        \item Language selector
        \item Perplexity-style sidebar with limited navigation
    \end{itemize}

    \item \textbf{Student Dashboard (Authenticated)}
    \begin{itemize}
        \item Feed-based recommendation display (Instagram-like scrolling)
        \item Upcoming registered events with calendar view
        \item Social feed from followed societies
        \item Sidebar navigation: Search, Explore, Societies, Events, Calendar
        \item Integrated calendar view showing registered events and events of interest
    \end{itemize}

    \item \textbf{Event Detail Page}
    \begin{itemize}
        \item Full event information with banner image
        \item Registration button with capacity indicator (login required)
        \item ``Add to Calendar'' button for calendar integration
        \item Society information card
        \item Similar events suggestions
    \end{itemize}

    \item \textbf{Calendar View}
    \begin{itemize}
        \item Monthly/weekly calendar displaying registered events
        \item Visual indicators for events with expressed interest
        \item ``Add to Personal Calendar'' integration (.ics export)
        \item Filter by society or event type
        \item Click event to view details
    \end{itemize}

    \item \textbf{Society Dashboard (Committee View)}
    \begin{itemize}
        \item Event/post management with CRUD actions
        \item Analytics overview with charts
        \item Follower metrics
        \item Additional ``Dashboard'' option in sidebar
    \end{itemize}

    \item \textbf{Social Feed}
    \begin{itemize}
        \item Card-based post layout with smooth scrolling
        \item Infinite scroll (engaging but not addictive)
        \item Like and ``Not Interested'' actions (authenticated only)
    \end{itemize}
\end{enumerate}

\textbf{Calendar Features Summary:}
\begin{itemize}
    \item In-platform calendar view for all registered events and events of interest
    \item Export to personal calendars via .ics file generation
    \item Compatible with Google Calendar, Apple Calendar, Outlook, and other iCalendar applications
    \item Automatic calendar entries for event registrations
    \item Optional calendar integration for expressed interest
\end{itemize}

\textit{Mockups to be created using Excalidraw/Figma and stored in \texttt{/mockups/} directory.}

\subsubsection{API Interface}

RESTful JSON API with JWT authentication.

\textbf{Endpoint Groups:}
\begin{itemize}
    \item \texttt{/api/auth/*} - OAuth login, callback, logout, token refresh
    \item \texttt{/api/users/*} - Profile management, interests
    \item \texttt{/api/events/*} - Event CRUD, registration, search
    \item \texttt{/api/posts/*} - Post CRUD, feed
    \item \texttt{/api/societies/*} - Society profiles, follow/unfollow, analytics
    \item \texttt{/api/recommendations} - Personalized content
    \item \texttt{/api/interactions} - Likes, ignores, feedback
    \item \texttt{/api/translate} - DeepL API proxy
\end{itemize}

\textbf{Rate Limiting:} 100 req/min (authenticated), 20 req/min (unauthenticated)

\textbf{Authentication:} JWT token in \texttt{Authorization: Bearer <token>} header

\subsubsection{External System Interfaces}

\begin{enumerate}
    \item \textbf{University OAuth (Microsoft Azure AD)}
    \begin{itemize}
        \item OAuth 2.0 authorization code flow
        \item Scopes: openid, profile, email, User.Read
        \item Callback URL configured in Azure portal
    \end{itemize}

    \item \textbf{Student Data API} (Subject to ISS approval)
    \begin{itemize}
        \item Sync student data: courses, department, year, society memberships
        \item Daily sync via cron job + manual refresh option
    \end{itemize}

    \item \textbf{DeepL Translation API}
    \begin{itemize}
        \item Auto-translate user-generated content
        \item English $\leftrightarrow$ Irish translations
        \item Usage monitoring to stay within quota
    \end{itemize}
\end{enumerate}

\subsubsection{Infrastructure}

\textbf{Current Deployment Architecture:}
\begin{itemize}
    \item \textbf{Hosting:} Self-hosted server accessible via \url{https://ct216.semyon.ie}
    \item \textbf{Cloudflare Tunnel:} Secure tunnel for external access without exposing server ports
    \item \textbf{Cloudflare CDN:} Global content delivery with automatic caching
    \item \textbf{Docker Containerization:} All services (Next.js, PostgreSQL, Redis) run in isolated containers
    \item \textbf{Network Security:} Dedicated subnet isolating databases from external access
    \item \textbf{SSL/TLS:} Cloudflare-managed certificates for HTTPS
\end{itemize}

\textbf{Recommendation Algorithm Implementation Strategy:}

\begin{itemize}
    \item \textbf{Initial Deployment:} Python-based implementation for rapid development and quick deployment
    \begin{itemize}
        \item Fast iteration and testing of recommendation logic
        \item Serves as production baseline during Semester 2 development
        \item Allows team to validate algorithm effectiveness early
    \end{itemize}

    \item \textbf{Optional C++/Rust Rewrite:} Personal development opportunity for project manager (Semyon)
    \begin{itemize}
        \item Timeline-dependent: only pursued if project timeline permits
        \item Primary purpose: systems programming language learning
        \item Not required for project completion or assessment
        \item Would provide performance improvements if implemented
    \end{itemize}
\end{itemize}

\textbf{Planned Infrastructure Improvements:}

\textit{Note: The following are preliminary plans and not definitive implementations. Final architecture may vary based on performance requirements and resource constraints.}

\begin{itemize}
    \item \textbf{Internal Nginx Proxy:}
    \begin{itemize}
        \item Nginx layer between Cloudflare and Next.js for enhanced caching
        \item Cache static assets locally to reduce Cloudflare CDN misses
        \item Load balancing across multiple Next.js instances if needed
        \item Additional request filtering and rate limiting
    \end{itemize}
\end{itemize}

\newpage
\section{Project Timeline}

\subsection{Project Overview}

\textbf{Total Duration:} Academic Year 2025/26 (Two Semesters)

\textbf{Semester 1 (Completed):}
\begin{itemize}
    \item Planning and requirements gathering
    \item Technology stack research and selection
    \item Infrastructure setup and deployment testing
    \item GitHub repository setup with CI/CD foundations
    \item Basic authentication system (bcrypt-based, custom implementation)
    \item PostgreSQL database (single table)
    \item Basic UI (home, login, register pages with generic styling)
    \item Deployment to production server (\url{https://ct216.semyon.ie})
\end{itemize}

\textbf{Semester 2 (12 Weeks):} Main development phase - all remaining features

\textbf{Final Submission:} End of Semester 2, Summer 2026

\subsection{Semester 2 Development Approach (12 Weeks)}

\textbf{Development Model:} Parallel feature development with 4 developers working independently on assigned features. Each developer owns 2-3 features end-to-end (frontend, backend, database).

\subsubsection{Infrastructure Setup Phase (Weeks 1-2)}
Shared infrastructure setup performed by team:
\begin{itemize}
    \item Database schema expansion and finalization
    \item OAuth migration from custom auth to Microsoft Azure AD
    \item Redis setup and caching infrastructure
    \item API rate limiting and authentication framework
\end{itemize}

\subsubsection{Parallel Feature Development (Weeks 3-10)}

\textbf{Feature Teams (Parallel):}

\begin{enumerate}
    \item \textbf{Events System} (Developer A)
    \begin{itemize}
        \item Event CRUD operations (create, read, update, delete)
        \item Registration and capacity management
        \item Event search and filtering
        \item Timeline: Weeks 3-10 (ongoing refinement)
    \end{itemize}

    \item \textbf{Societies \& Social Posts} (Developer B)
    \begin{itemize}
        \item Society profiles and management
        \item Follow/unfollow functionality
        \item Post creation and feed display
        \item Post interactions (like, ignore)
        \item Timeline: Weeks 3-10 (ongoing refinement)
    \end{itemize}

    \item \textbf{Recommendation Engine} (Developer C)
    \begin{itemize}
        \item Python-based algorithm implementation
        \item Batch processing and caching with Redis
        \item Scoring logic based on interests, attendance, department
        \item Timeline: Weeks 3-10 (ongoing refinement and optimization)
    \end{itemize}

    \item \textbf{Multilingual Support \& Calendar Integration} (Developer D)
    \begin{itemize}
        \item next-intl configuration and UI language switching
        \item DeepL API integration for auto-translation
        \item Calendar view and .ics export functionality
        \item UI/UX polish and accessibility improvements
        \item Timeline: Weeks 3-10 (ongoing refinement)
    \end{itemize}
\end{enumerate}

\textbf{Integration Points:}
\begin{itemize}
    \item Weekly code reviews and integration testing
    \item API contract testing between features
    \item Shared database and Redis access
    \item Continuous integration via GitHub Actions
\end{itemize}

\subsection{Key Milestones}

\begin{table}[h]
\centering
\begin{tabular}{|c|c|p{8cm}|}
\hline
\textbf{Milestone} & \textbf{Week} & \textbf{Deliverable} \\
\hline
M0 (Complete) & S1 & Basic auth, database, deployment, planning \\
\hline
M1 & Week 2 & Infrastructure setup complete: OAuth, Redis, database schema finalized \\
\hline
M2 & Week 6 & All features in alpha: Events, Societies/Posts, Recommendations, Multilingual \\
\hline
M3 & Week 9 & All features functionally complete with integration testing \\
\hline
M4 & Week 10 & UI/UX polish, accessibility, performance optimization \\
\hline
M5 & Week 11 & Comprehensive testing, bug fixes, production readiness \\
\hline
M6 & Week 12 & Final documentation, demo, submission \\
\hline
\end{tabular}
\end{table}

\subsection{Gantt Charts}

\textit{[PLACEHOLDER: Detailed Gantt charts to be created and maintained separately]}

\subsection{Risk Management}

\begin{table}[h]
\centering
\begin{tabular}{|p{4cm}|c|c|p{5cm}|}
\hline
\textbf{Risk} & \textbf{Prob.} & \textbf{Impact} & \textbf{Mitigation} \\
\hline
University OAuth approval delayed & Medium & High & Continue with custom auth, add OAuth if approved \\
\hline
Student API access denied & High & Medium & Manual profile setup, no automatic sync \\
\hline
Timeline slippage (12 weeks) & High & High & Prioritize core features: events, societies, basic recommendations \\
\hline
Team member unavailable & Medium & Medium & PM can assist development, task redistribution \\
\hline
Recommendation algorithm complexity & Medium & High & Start with simple scoring, optimize later \\
\hline
Multilingual implementation time & Medium & Medium & English-only MVP, add Irish if time permits \\
\hline
\end{tabular}
\end{table}

\textbf{MVP Core Features (Must-Have):}
\begin{itemize}
    \item Events system (CRUD, registration, basic search)
    \item Societies (profiles, following)
    \item Basic recommendation algorithm
    \item Responsive UI with feed layout
\end{itemize}

\textbf{Secondary Features (Nice-to-Have):}
\begin{itemize}
    \item Social posts and interactions
    \item Calendar integration (.ics export)
    \item Analytics dashboard
    \item Multilingual support (Irish)
    \item DeepL auto-translation
    \item University OAuth migration
\end{itemize}

\textbf{Stretch Goals (If Time Permits):}
\begin{itemize}
    \item Advanced recommendation algorithm (Python/Rust rewrite)
    \item Internal Nginx caching layer
    \item Push notifications
    \item Student API integration
\end{itemize}

\newpage
\section*{A. Appendices}

\textit{The information contained in these appendices is supplementary and informational. It is not considered part of the formal requirements unless explicitly referenced in Sections 1-3.}

\subsection*{A.1 Glossary}

\begin{table}[h]
\centering
\begin{tabular}{|p{3cm}|p{10cm}|}
\hline
\textbf{Term} & \textbf{Definition} \\
\hline
API & Application Programming Interface \\
\hline
CSRF & Cross-Site Request Forgery - security vulnerability \\
\hline
GDPR & General Data Protection Regulation - EU data privacy law \\
\hline
JWT & JSON Web Token - authentication token format \\
\hline
OAuth & Open Authorization - delegated authentication standard \\
\hline
RBAC & Role-Based Access Control \\
\hline
Redis & In-memory data store for caching and sessions \\
\hline
SSR/CSR & Server-Side Rendering / Client-Side Rendering \\
\hline
TTL & Time To Live - cache expiration time \\
\hline
\end{tabular}
\end{table}

\subsection*{A.2 Technology Stack Summary}

\begin{table}[h]
\centering
\begin{tabular}{|p{3cm}|p{10cm}|}
\hline
\textbf{Layer} & \textbf{Technology} \\
\hline
Frontend & Next.js 14+, React 18, Bootstrap 5, TypeScript \\
\hline
Backend & Next.js API Routes, JWT \\
\hline
Database & PostgreSQL 15+ \\
\hline
Cache & Redis 7+ \\
\hline
Proxy & Nginx \\
\hline
DevOps & Docker, GitHub Actions \\
\hline
Cloud & AWS (ECS/RDS/ElastiCache) \\
\hline
\end{tabular}
\end{table}

\subsection*{A.3 Database Tables Overview}

Core tables (full schema in ARCHITECTURE\_AND\_TECH\_STACK.md):

\begin{itemize}
    \item \texttt{users} - Student/member profiles
    \item \texttt{student\_data} - Synced university data
    \item \texttt{societies} - Society information
    \item \texttt{society\_members} - Memberships and roles
    \item \texttt{events} - Event metadata
    \item \texttt{event\_translations} - Multilingual event content
    \item \texttt{posts} - Social posts with JSONB translations
    \item \texttt{registrations} - Event attendance
    \item \texttt{interactions} - Likes, ignores, interests
    \item \texttt{recommendations} - Pre-computed suggestions
    \item \texttt{user\_interests} - User-selected interests
\end{itemize}

\end{document}
